\begin{abstract}
Choosing between classical and Bayesian sparse regression methods involves a real trade-off: penalized estimators like Lasso run in milliseconds but give no uncertainty estimates, while Horseshoe and Spike-and-Slab priors produce full posteriors but need MCMC chains that take minutes per fit.
Surprisingly few studies compare these two families head-to-head under the conditions that actually make sparse regression hard---correlated features, weak signals, and growing dimensionality.
We benchmark six methods (OLS, Ridge, Lasso, Elastic Net, Horseshoe, Spike-and-Slab) on synthetic data with three covariance structures ($\rho$ up to 0.9), four SNR levels, and $p \in \{20, 50, 100\}$, plus the Diabetes dataset, totalling over 2{,}600 experiments.
The results are clear on some points and nuanced on others.
Bayesian methods win on prediction error (MSE 72 vs.\ 108--267), and the Horseshoe delivers near-nominal 95\% coverage (94.8\%).
But Spike-and-Slab, despite narrower intervals, under-covers at 91.9\%---its continuous relaxation likely plays a role.
For variable selection, Lasso and Spike-and-Slab tie at $F_1 \approx 0.47$, making Lasso the practical default when posteriors are not needed.
Code and data are available at \url{https://github.com/xiao98/sparse-bayesian-regression-bench}.
\end{abstract}
